\documentclass[11pt,a4paper]{article}

\usepackage[utf8]{inputenc}
\usepackage[T1]{fontenc}
\usepackage[spanish,es-noshorthands,es-nodecimaldot]{babel}
\usepackage[margin=2.5cm]{geometry}
\usepackage{graphicx}
\usepackage{float}
\usepackage{booktabs}
\usepackage{amsmath}
\usepackage{enumitem}
\usepackage{caption}
\usepackage[hidelinks]{hyperref}

% Las figuras y tablas se generan en ../output/
\graphicspath{{../output/}}

\setlist[itemize]{itemsep=2pt,topsep=4pt}
\captionsetup{font=small,labelfont=bf}

\begin{document}

\begin{titlepage}
\centering
\vspace*{3cm}

{\LARGE\bfseries Trabajo Práctico N$^{\circ}$3\par}
\vspace{0.4cm}
{\large (Semestre de Otoño 2025)\par}
\vspace{0.3cm}
{\large\bfseries Big Data\par}

\vfill

{\large Profesora: María Noelia Romero\par}
\vspace{0.2cm}
{\large Tutor: Tomás Pacheco\par}
\vspace{0.2cm}
{\large Alumnos: Catalina Gallardo, Manuel Elorga\par}

\vfill
\end{titlepage}

% ===============================================================
\section{Un poco más de estadística descriptiva}
% ===============================================================

\subsection*{Ejercicio 1}

\begin{itemize}
  \item \textbf{Sexo (CH04)}: indica el sexo del individuo. Es útil para considerar posibles
        disparidades de género en el acceso y la permanencia en el mercado laboral.

  \item \textbf{Edad (CH06)}: registra la edad del individuo. Es relevante para distinguir
        entre población activa e inactiva y para detectar patrones etarios en la desocupación.

  \item \textbf{Estado civil (CH07)}: refleja la situación conyugal del individuo. Puede ser
        relevante si la desocupación afecta de forma distinta a personas solteras, casadas,
        divorciadas o viudas.

  \item \textbf{Lugar de residencia hace cinco años (CH16)}: informa dónde vivía la persona
        cinco años antes de la encuesta. Permite identificar movilidad geográfica, dentro del
        país o desde el exterior, lo cual puede incidir en las posibilidades de inserción
        laboral, el acceso a redes de empleo y la adaptación al mercado de trabajo local.

  \item \textbf{Cantidad de personas en el hogar (PERSONAS\_HOGAR)}: indica cuántas personas
        viven en el hogar del encuestado. Puede estar relacionada con la presión económica del
        hogar y con la necesidad de que sus miembros participen del mercado laboral. Como se
        ve en la Figura~\ref{fig:personas}, se presenta una distribución sesgada hacia la
        derecha, en la que los hogares tienen entre dos y seis miembros. Sin embargo, existen
        valores extremos que aumentan la varianza.

  \item \textbf{Edad al cuadrado (EDAD2)}: corresponde al cuadrado de la edad del individuo.
        Permite capturar posibles efectos no lineales de la edad sobre la condición de
        ocupación, como diferencias entre jóvenes, adultos y personas próximas a la
        jubilación. Como se ve en la Figura~\ref{fig:edad2}, hay una curtosis alta y una cola
        larga, lo que confirma que la muestra es relativamente joven.

  \item \textbf{Obra social (OBRA\_SOCIAL)}: variable binaria que toma valor 0 si la persona
        depende total o parcialmente de los servicios públicos de salud. Puede ser un buen
        predictor de ocupación, ya que contar con obra social suele estar vinculado a una
        relación laboral formal. En la Figura~\ref{fig:obra} se visualiza que aproximadamente
        el 60\% de los individuos carece de obra social privada.

  \item \textbf{Años de educación (EDUCACION)}: estima la cantidad de años de educación
        formal del individuo. Fue ajustada para reducir valores perdidos. Esta variable es
        clave, ya que el nivel educativo influye fuertemente en las oportunidades laborales y
        en la probabilidad de inserción en el mercado formal. En la
        Figura~\ref{fig:educacion} se puede ver un patrón en las personas que tienen la
        primaria completa, la secundaria completa y estudios superiores, lo que evidencia la
        heterogeneidad en el capital humano.

  \item \textbf{Interacción entre edad y educación (EDADxEDUCACION)}: representa el producto
        entre la edad y los años de educación. Busca capturar efectos combinados, como que
        una mayor edad acompañada de mayor nivel educativo se traduzca en más experiencia y,
        por lo tanto, en una mayor probabilidad de empleo. La Figura~\ref{fig:edadxeduc}
        muestra que la interacción entre los valores medianos de los ocupados es más alta que
        la de los desocupados, lo que sugiere que una mayor formación y experiencia
        incrementan la empleabilidad.

  \item \textbf{Educación en colegio privado (COLEGIO\_PRIVADO)}: variable binaria que indica
        si el individuo asistió a una institución educativa privada. Puede utilizarse como
        \textit{proxy} del nivel socioeconómico del hogar de origen, lo que incide en las
        trayectorias educativas y laborales futuras. En la Figura~\ref{fig:privado} vemos que
        la gran mayoría asistió a una institución educativa pública: solo el 10\% de la
        muestra reporta haber asistido a un colegio privado.
\end{itemize}

\begin{figure}[H]
  \centering
  \begin{minipage}{0.48\textwidth}
    \centering
    \includegraphics[width=\textwidth]{{personas-hogar}.png}
    \caption{Cantidad de personas en el hogar.}
    \label{fig:personas}
  \end{minipage}\hfill
  \begin{minipage}{0.48\textwidth}
    \centering
    \includegraphics[width=\textwidth]{{edad2}.png}
    \caption{Distribución de la edad al cuadrado.}
    \label{fig:edad2}
  \end{minipage}

  \vspace{0.6cm}

  \begin{minipage}{0.48\textwidth}
    \centering
    \includegraphics[width=\textwidth]{{obra-social}.png}
    \caption{Cobertura de obra social.}
    \label{fig:obra}
  \end{minipage}\hfill
  \begin{minipage}{0.48\textwidth}
    \centering
    \includegraphics[width=\textwidth]{{educacion}.png}
    \caption{Años de educación.}
    \label{fig:educacion}
  \end{minipage}

  \vspace{0.6cm}

  \begin{minipage}{0.48\textwidth}
    \centering
    \includegraphics[width=\textwidth]{{edadxeducacion}.png}
    \caption{Interacción entre edad y educación.}
    \label{fig:edadxeduc}
  \end{minipage}\hfill
  \begin{minipage}{0.48\textwidth}
    \centering
    \includegraphics[width=\textwidth]{{colegio-privado}.png}
    \caption{Asistencia a colegio privado.}
    \label{fig:privado}
  \end{minipage}
\end{figure}

% ===============================================================
\section{Enfoque de validación}
% ===============================================================

\subsection*{Ejercicio 2}

El Cuadro~\ref{tab:medias} presenta un análisis de las diferencias de medias entre la base de
entrenamiento (\textit{train}) y la base de prueba (\textit{test}) para las variables
elegidas, en 2004 y 2024. El objetivo es verificar que la partición de los datos realizada
con \texttt{train\_test\_split} fue aleatoria y que ambas muestras son comparables.

Las diferencias de medias entre \textit{train} y \textit{test} son en general muy pequeñas,
lo cual indica un buen balance. Por ejemplo, variables como CH04 (sexo), CH07 (estado civil),
PERSONAS\_HOGAR, OBRA\_SOCIAL y COLEGIO\_PRIVADO presentan diferencias nulas o cercanas a
cero, lo cual es esperable en una partición aleatoria.

En 2004, algunas variables muestran diferencias algo mayores, como EDAD2 ($-54.58$) y
EDUCACION ($0.41$). No obstante, en el caso de EDAD2 esto puede explicarse por una pequeña
diferencia en la edad promedio, amplificada por la transformación cuadrática. En cuanto a
EDUCACION, si bien la diferencia es la más elevada en términos absolutos, sigue siendo pequeña
y no se considera problemática.

En 2024 las diferencias son aún menores. La mayor corresponde a CH06 (edad), con $-0.44$,
mientras que EDAD2 y EDUCACION presentan diferencias de $-11.72$ y $-0.09$ respectivamente,
todas dentro de márgenes aceptables. Variables como EDADxEDUCACION ($-4.06$) también muestran
ligeras diferencias, pero no significativas.

En conclusión, todas las diferencias observadas son pequeñas y no estadísticamente
significativas. Esto confirma que la división entre las bases de entrenamiento y prueba fue
adecuada y que ambas muestras son representativas del total. Por lo tanto, pueden ser
utilizadas sin inconvenientes en el desarrollo y la validación de modelos predictivos de la
condición de ocupación.

\begin{table}[H]
  \centering
  \caption{Medias y diferencias entre las muestras de entrenamiento y testeo, 2004 y 2024.}
  \label{tab:medias}
  \small
  \begin{tabular}{lrrrrrr}
    \toprule
    & \multicolumn{3}{c}{\textbf{2004}} & \multicolumn{3}{c}{\textbf{2024}} \\
    \cmidrule(lr){2-4}\cmidrule(lr){5-7}
    Variable & Media train & Media test & Diferencia & Media train & Media test & Diferencia \\
    \midrule
    CONSTANTE        &    1.00 &    1.00 &   0.00 &    1.00 &    1.00 &   0.00 \\
    SEXO             &    1.52 &    1.52 &   0.00 &    1.52 &    1.53 &  -0.01 \\
    EDAD             &   31.29 &   31.61 &  -0.32 &   36.22 &   36.66 &  -0.44 \\
    ESTADO CIVIL     &    3.67 &    3.63 &   0.04 &    3.55 &    3.51 &   0.04 \\
    DONDE VIVE       &    1.47 &    1.50 &  -0.03 &    1.31 &    1.26 &   0.05 \\
    PERSONAS HOGAR   &    4.83 &    4.88 &  -0.04 &    4.06 &    4.11 &  -0.05 \\
    EDAD2            & 1430.17 & 1484.76 & -54.58 & 1834.88 & 1846.60 & -11.72 \\
    OBRA SOCIAL      &    0.56 &    0.56 &  -0.00 &    0.65 &    0.64 &   0.01 \\
    EDUCACION        &    7.92 &    7.51 &   0.41 &    9.06 &    9.15 &  -0.09 \\
    EDADxEDUCACION   &  279.93 &  273.28 &   6.64 &  367.69 &  371.75 &  -4.06 \\
    COLEGIO PRIVADO  &    0.06 &    0.05 &   0.01 &    0.06 &    0.06 &   0.00 \\
    \bottomrule
  \end{tabular}
\end{table}

% ===============================================================
\section{Métodos de clasificación y performance}
% ===============================================================

\subsection*{Ejercicio 3}

Para 2004, el Cuadro~\ref{tab:metricas04} reporta el \textit{accuracy} y el AUC de cada
modelo, el Cuadro~\ref{tab:cm04} las matrices de confusión y la Figura~\ref{fig:roc04} las
curvas ROC.

\begin{table}[H]
  \centering
  \caption{Métricas de desempeño, 2004.}
  \label{tab:metricas04}
  \begin{tabular}{lrr}
    \toprule
    Modelo & Accuracy & AUC \\
    \midrule
    Regresión logística                      & 0.9544 & 0.7631 \\
    Análisis discriminante lineal (LDA)      & 0.9544 & 0.7643 \\
    Análisis discriminante cuadrático (QDA)  & 0.7630 & 0.7508 \\
    $K$-vecinos más cercanos ($K=5$)         & 0.9509 & 0.5829 \\
    Naive Bayes                              & 0.8252 & 0.7178 \\
    \bottomrule
  \end{tabular}
\end{table}

\begin{table}[H]
  \centering
  \caption{Matrices de confusión, 2004.}
  \label{tab:cm04}
  \begin{tabular}{lcccc}
    \toprule
    & \multicolumn{2}{c}{Real 0} & \multicolumn{2}{c}{Real 1} \\
    \cmidrule(lr){2-3}\cmidrule(lr){4-5}
    Modelo & Pred 0 & Pred 1 & Pred 0 & Pred 1 \\
    \midrule
    Regresión logística              & 1381 &   0 & 66 &  0 \\
    LDA                              & 1381 &   0 & 66 &  0 \\
    QDA                              & 1065 & 316 & 27 & 39 \\
    $K$-vecinos ($K=5$)              & 1375 &   6 & 65 &  1 \\
    Naive Bayes                      & 1162 & 219 & 34 & 32 \\
    \bottomrule
  \end{tabular}
\end{table}

\begin{figure}[H]
  \centering
  \begin{minipage}{0.32\textwidth}
    \centering
    \includegraphics[width=\textwidth]{{reg-log-2004}.png}\\[-2pt]
    {\scriptsize Regresión logística}
  \end{minipage}\hfill
  \begin{minipage}{0.32\textwidth}
    \centering
    \includegraphics[width=\textwidth]{{LDA-2004}.png}\\[-2pt]
    {\scriptsize LDA}
  \end{minipage}\hfill
  \begin{minipage}{0.32\textwidth}
    \centering
    \includegraphics[width=\textwidth]{{QDA-2004}.png}\\[-2pt]
    {\scriptsize QDA}
  \end{minipage}

  \vspace{0.4cm}

  \begin{minipage}{0.32\textwidth}
    \centering
    \includegraphics[width=\textwidth]{{KNN-2004}.png}\\[-2pt]
    {\scriptsize $K$-vecinos ($K=5$)}
  \end{minipage}\hfill
  \begin{minipage}{0.32\textwidth}
    \centering
    \includegraphics[width=\textwidth]{{NB-2004}.png}\\[-2pt]
    {\scriptsize Naive Bayes}
  \end{minipage}\hfill
  \begin{minipage}{0.32\textwidth}~\end{minipage}

  \caption{Curvas ROC, 2004.}
  \label{fig:roc04}
\end{figure}

Para 2024, el Cuadro~\ref{tab:metricas24} reporta las métricas, el Cuadro~\ref{tab:cm24} las
matrices de confusión y la Figura~\ref{fig:roc24} las curvas ROC.

\begin{table}[H]
  \centering
  \caption{Métricas de desempeño, 2024.}
  \label{tab:metricas24}
  \begin{tabular}{lrr}
    \toprule
    Modelo & Accuracy & AUC \\
    \midrule
    Regresión logística                      & 0.9815 & 0.8677 \\
    Análisis discriminante lineal (LDA)      & 0.9815 & 0.8670 \\
    Análisis discriminante cuadrático (QDA)  & 0.8234 & 0.8121 \\
    $K$-vecinos más cercanos ($K=5$)         & 0.9815 & 0.5262 \\
    Naive Bayes                              & 0.8647 & 0.7775 \\
    \bottomrule
  \end{tabular}
\end{table}

\begin{table}[H]
  \centering
  \caption{Matrices de confusión, 2024.}
  \label{tab:cm24}
  \begin{tabular}{lcccc}
    \toprule
    & \multicolumn{2}{c}{Real 0} & \multicolumn{2}{c}{Real 1} \\
    \cmidrule(lr){2-3}\cmidrule(lr){4-5}
    Modelo & Pred 0 & Pred 1 & Pred 0 & Pred 1 \\
    \midrule
    Regresión logística              & 1328 &   0 & 25 &  0 \\
    LDA                              & 1328 &   0 & 25 &  0 \\
    QDA                              & 1096 & 232 &  7 & 18 \\
    $K$-vecinos ($K=5$)              & 1328 &   0 & 25 &  0 \\
    Naive Bayes                      & 1155 & 173 & 10 & 15 \\
    \bottomrule
  \end{tabular}
\end{table}

\begin{figure}[H]
  \centering
  \begin{minipage}{0.32\textwidth}
    \centering
    \includegraphics[width=\textwidth]{{reg-log-2024}.png}\\[-2pt]
    {\scriptsize Regresión logística}
  \end{minipage}\hfill
  \begin{minipage}{0.32\textwidth}
    \centering
    \includegraphics[width=\textwidth]{{LDA-2024}.png}\\[-2pt]
    {\scriptsize LDA}
  \end{minipage}\hfill
  \begin{minipage}{0.32\textwidth}
    \centering
    \includegraphics[width=\textwidth]{{QDA-2024}.png}\\[-2pt]
    {\scriptsize QDA}
  \end{minipage}

  \vspace{0.4cm}

  \begin{minipage}{0.32\textwidth}
    \centering
    \includegraphics[width=\textwidth]{{KNN-2024}.png}\\[-2pt]
    {\scriptsize $K$-vecinos ($K=5$)}
  \end{minipage}\hfill
  \begin{minipage}{0.32\textwidth}
    \centering
    \includegraphics[width=\textwidth]{{NB-2024}.png}\\[-2pt]
    {\scriptsize Naive Bayes}
  \end{minipage}\hfill
  \begin{minipage}{0.32\textwidth}~\end{minipage}

  \caption{Curvas ROC, 2024.}
  \label{fig:roc24}
\end{figure}

\subsection*{Ejercicio 4}

Para determinar qué método predice mejor en cada año es importante analizar varias métricas
que reflejan el desempeño del modelo, sobre todo considerando que la clase desocupado es
minoritaria y los datos están desbalanceados. Las métricas más relevantes para esta tarea son
el \textit{recall}, la precisión, el F1-score, el área bajo la curva ROC y la matriz de
confusión. Otra medida relevante en este sentido es el error cuadrático medio (MSE), que
indica los errores cometidos en la predicción.

En 2004, los métodos de regresión logística y análisis discriminante lineal tienen un
\textit{accuracy} alto (aproximadamente 95\%) y alta precisión, pero un \textit{recall} de 0
para la clase desocupado, ya que nunca predicen 1. El método de vecinos más cercanos también
presenta un \textit{recall} muy bajo, casi nulo. El MSE más bajo entre los modelos corresponde
a la regresión logística; sin embargo, esto se debe a que el modelo nunca predice 1, por lo
que sus únicos errores son todos los desocupados, a los que asigna la categoría de no
desocupados.

Por otro lado, Naive Bayes muestra un \textit{recall} de 0.485 y un F1-score de 0.202, lo que
refleja un mejor equilibrio entre detectar desocupados y evitar falsos positivos que los
métodos anteriores. Sin embargo, el método que logra el mejor \textit{recall} es QDA, con
0.591: identifica un mayor porcentaje de desocupados, aunque su precisión es baja y el
F1-score es moderado, lo que indica que comete más errores clasificando ocupados como
desocupados. En conclusión, para 2004 el equilibrio entre precisión y \textit{recall} se
alcanza con Naive Bayes; pero como lo que buscamos es predecir desocupados, QDA parecería ser
la mejor opción, aunque se cometerán errores de sobrestimación de la desocupación.

En 2024, nuevamente la regresión logística, LDA y KNN presentan un \textit{accuracy} muy alto
(cerca del 98\%) pero un \textit{recall} igual a 0 para la clase desocupado, por lo que nunca
predicen 1. Naive Bayes mejora el \textit{recall} a 0.6 con un F1-score de 0.141, lo que
muestra una mejor capacidad para detectar desocupados que los tres anteriores. De nuevo, la
regresión logística tiene el menor MSE. Al igual que en 2004, QDA tiene el \textit{recall} más
alto, de 0.72, con lo que identifica el 72\% de los desocupados, aunque con un F1-score
ligeramente menor debido a su baja precisión. Por lo tanto, en 2024 QDA es el método que mejor
detecta la clase desocupado debido a su alto \textit{recall}, mientras que Naive Bayes ofrece
un mejor balance entre precisión y \textit{recall}.

En problemas con clases desbalanceadas, como este, el \textit{recall} y el F1-score son
métricas más informativas que el \textit{accuracy}, ya que este último puede ser engañoso si
el modelo predice principalmente la clase mayoritaria.

\subsection*{Ejercicio 5}

Al no tener observaciones de la base \texttt{norespondieron} en 2004, predecimos únicamente en
2024. Vamos a utilizar dos métodos: QDA, que es el mejor prediciendo desocupados pero
sobrestima, y la regresión logística, que es la que menor error cuadrático medio tiene pero
subestima.

El modelo QDA predice que, de los que no respondieron, el 50.0\% serán desocupados. Por otro
lado, el modelo de regresión logística predice que, de los que no respondieron, el 0\% serán
desocupados. Vemos que en esta muestra también se cumple que la regresión logística no predice
ningún desocupado y que QDA, probablemente, sobrestima. Es importante considerar que estas
estimaciones no son muy representativas, ya que la muestra tiene solo seis observaciones.

\end{document}
